\section*{第 \ref{ch:b3:origin-of-elements} 章——元素从何而来}
% Values from the Book 5 scratch calculator calc/ch30.py.

\begin{solution}{exo:b3:origin-of-elements:1}
(a)~$\mathrm{Si/Fe} = 10^{7.51 - 7.46} = 10^{0.05} = 1.12$ 和
$\mathrm{O/Fe} = 10^{8.69 - 7.46} = 10^{1.23} = 17.0$：每个铁原子约对应一个硅
原子和十七个氧原子。(b)~$\mathrm{Fe/Au} =
10^{7.46 - 0.91} = 10^{6.55} = \num{3.55e6}$ 和 $\mathrm{Pb/Au} =
10^{1.95 - 0.91} = 10^{1.04} = 11.0$：超过铁峰后，丰度
骤降数百万倍，但铅位于其自身的峰上，丰度
是金的十一倍。
\end{solution}

\begin{solution}{exo:b3:origin-of-elements:2}
(a)~$10^{7.46 - 12} = \num{2.884e-5}$，即每百万个氢原子中
有 \num{28.8} 个铁原子。(b)~少一百倍，即每百万个氢原子中仅有 \num{0.288} 个铁原子，
约每 $1/\num{2.884e-7} = \num{3.47e6}$ 个氢原子中
有一个铁原子。(c)~$[\mathrm{Fe/H}] = \log_{10}(\num{e-7}/\num{2.884e-5}) =
-2.46$。
\end{solution}

\begin{solution}{exo:b3:origin-of-elements:3}
$^{124}\mathrm{Xe}$ 和 $^{126}\mathrm{Xe}$ 位于富质子一侧，
不在 \omterm{def:b3:origin-of-elements:s-process}{s 过程}路径上，并被稳定的 $^{124}\mathrm{Te}$ 和
$^{126}\mathrm{Te}$：\omterm{def:b3:origin-of-elements:p-process}{p 核}。$^{128}\mathrm{Xe}$ 和 $^{130}\mathrm{Xe}$
位于 \omterm{def:b3:origin-of-elements:s-process}{s 过程}路径上，但被稳定的
$^{128}\mathrm{Te}$ 和 $^{130}\mathrm{Te}$ 屏蔽于 \omterm{def:b3:origin-of-elements:s-process}{r 过程}之外：仅 \omterm{def:b3:origin-of-elements:s-process}{s 过程}。$^{129}\mathrm{Xe}$、
$^{131}\mathrm{Xe}$ 和 $^{132}\mathrm{Xe}$ 位于 \omterm{def:b3:origin-of-elements:s-process}{s 过程}路径上，且没有
较低 $Z$ 的稳定同量异位素（碘只有 127）：两种过程均参与。
$^{134}\mathrm{Xe}$ 和 $^{136}\mathrm{Xe}$ 位于
$^{133}\mathrm{Xe}$ 之外（\omterm{def:b3:origin-of-elements:s-process}{s 过程}路径在此离开氙），它们未被
屏蔽：仅 \omterm{def:b3:origin-of-elements:s-process}{r 过程}。
\end{solution}

\begin{solution}{exo:b3:origin-of-elements:4}
(a)~$t = d/c = 40\times\qty{3.086e22}{m}/\qty{2.998e8}{m.s^{-1}} =
\qty{4.12e15}{s}$，约 1.3 亿年。(b)~$\Delta v/c \le
1.7/\num{4.12e15} = \num{4.1e-16}$。这是一个上限，因为部分
或全部延迟可能来自伽马射线在
并合之后的发射。
\end{solution}

\begin{solution}{exo:b3:origin-of-elements:5}
(a)~$295\times0.1067 = \qty{31.48}{mb}$，$832\times0.0373 =
\qty{31.03}{mb}$，$155\times0.1984 = \qty{30.75}{mb}$。(b)~
由 $^{122}\mathrm{Te}$ 出发：$N_{123} = 31.48/832 = \num{0.03783}$，\qty{1.4}{\percent}
高于实测值；$N_{124} = 31.48/155 = \num{0.2031}$，
高出 \qty{2.4}{\percent}。仅由截面即可将相差五倍的丰度
预测精度可达几个百分点以内。
(c)~$^{128}\mathrm{Te}$ 仅由 \omterm{def:b3:origin-of-elements:s-process}{r 过程}形成：s 流在碲
$^{127}\mathrm{Te}$ 处离开，该核衰变，因此没有稳定的 s 流经过
$^{128}\mathrm{Te}$，其丰度由 \omterm{def:b3:origin-of-elements:s-process}{r 过程}决定。
\end{solution}

\begin{solution}{exo:b3:origin-of-elements:6}
(a)~$\tau_0\langle\sigma\rangle = 0.306\times155 = 47.4$，故
$f = 1/(1 + 1/47.4) = \num{0.979}$，对于 $^{124}\mathrm{Te}$；对
$^{138}\mathrm{Ba}$，$0.306\times4.0 = 1.224$，则 $f = 1/(1 + 0.817) =
\num{0.550}$：幻数核仅传递 \qty{55}{\percent} 的
流。(b)~六个因子分别为 \num{0.959}、\num{0.550}、\num{0.908}、
\num{0.771}、\num{0.972} 和 \num{0.915}；其乘积为 \num{0.328}：
$\sigma N$ 下降了 \num{3.05} 倍。(c)~实测下降倍数为
$27.6/7.72 = 3.57$。模型对台阶大小的估计误差在
\qty{15}{\percent} 之内，其中几乎全部来自 $^{138}\mathrm{Ba}$ 和
$^{140}\mathrm{Ce}$，这两个中子数 $N = 82$ 的核。
\end{solution}

\begin{solution}{exo:b3:origin-of-elements:7}
(a)~$kT = 8\times10^3\times\qty{1.602e-19}{J} = \qty{1.282e-15}{J}$，并
$v_T = \sqrt{2\times\num{1.282e-15}/\num{1.6749e-27}} =
\qty{1.24e6}{m.s^{-1}}$。(b)~取 $\langle\sigma\rangle =
\qty{155e-31}{m^2}$，$\lambda = n_\mathrm{n}\langle\sigma\rangle v_T$ 为
\qty{1.92e-11}{s^{-1}}，当 $n_\mathrm{n} = \qty{e12}{m^{-3}}$
（$10^6$ 个每立方厘米）：平均时间为 \qty{5.2e10}{s}，约
\num{1650} 年；随后在 $10^8$ 时为 \num{16.5} 年，在 \num{0.165}~年，
两个月，在 \qty{e10}{cm^{-3}} 时。(c)~$^{127}\mathrm{Te}$ 在
小时内衰变，远快于这些俘获时间：它总是先衰变，
因此路径在此不分叉。$^{134}\mathrm{Cs}$ 的寿命约为两
年，介于 $10^8$ 和 \qty{e10}{cm^{-3}} 时的俘获时间之间：在
氦闪的高密度下，其中相当一部分在
衰变之前被俘获。它是一个分支点，其两侧的丰度
可用来测定中子密度。
\end{solution}

\begin{solution}{exo:b3:origin-of-elements:8}
(a)~比钙稀少 $10^{6.30 - 3.14} = 1445$ 倍，并
比钛稀少 $10^{4.97 - 3.14} = 68$ 倍。(b)~$10^{8.46 - 1.38}
= \num{1.2e7}$。(c)~钪显示奇偶效应：质子数为
偶数的核结合更紧密，由大质量恒星燃烧阶段的
氦核俘获形成，这些阶段有利于它们。铍
呈现锂、铍和硼的低谷：
这些脆弱的核在恒星
内部温度下被破坏，氦燃烧直接由氦变为碳，
不会产生它们。
\end{solution}

\begin{solution}{exo:b3:origin-of-elements:9}
(a)~设 $\alpha$ 为 \omterm{def:b3:origin-of-elements:alpha}{α 元素}，$\mathrm{Fe}_\mathrm{cc}$
和 $\mathrm{Fe}_\mathrm{Ia}$ 分别为两类超新星产生的铁。在
平台段上，$\alpha/\mathrm{Fe}_\mathrm{cc} = 10^{P}$ 乘以
太阳比值。在 $[\mathrm{Fe/H}] = 0$ 处，$\alpha/(\mathrm{Fe}_\mathrm{cc} +
\mathrm{Fe}_\mathrm{Ia})$ 为太阳值。两式相除，
$\mathrm{Fe}_\mathrm{cc}/(\mathrm{Fe}_\mathrm{cc} + \mathrm{Fe}_\mathrm{Ia})
= 10^{-P}$，且 Ia 型超新星所占份额为 $1 - 10^{-P} = 1 - 10^{-0.4} =
\num{0.602}$。(b)~$10^{-P} = 1/3$ 得 $P = \log_{10}3 = 0.48$。
(c)~两个星系中的 Ia 型超新星均在相同的延迟之后开始出现。
当它们在矮星系中开始出现时，该星系的气体仅
达到 $[\mathrm{Fe/H}] = -2$，而银河系中为 $-1$：该星系
把气体转化为恒星的速度慢了十倍。
\end{solution}

\begin{solution}{exo:b3:origin-of-elements:10}
(a)~对于锆，
\[
\frac{(X/\mathrm{Mg})_\mathrm{Earth}}{(X/\mathrm{Mg})_\mathrm{CI}} =
\frac{10.5/228000}{3.96/95870} = 1.11 ;
\]
类似地，钠为 $(2670/228000)/(5010/95870) = 0.224$，
钾为 \num{0.190}，铷为 \num{0.118}。(b)~锆是最难熔的，其比值
接近球粒陨石的比值。三种中等挥发性的
碱金属在 \qtyrange{800}{1006}{K} 下凝结，其丰度相对于球粒陨石
降低了五至九倍，铷是挥发性最强的，降低得最多。
硅酸盐地球遵循\omterm{def:b3:origin-of-elements:volatility}{挥发性趋势}。(c)~镁是主要的难熔元素，
与岩石主体一同凝结，且它是
亲石元素：地核形成时，它留在地幔中。若以
铁或镍之类的亲铁元素为参照，会把流入地核的损失
与挥发性造成的损失混在一起。
\end{solution}

\begin{solution}{exo:b3:origin-of-elements:11}
(a)~当质量 $\mathrm{d}s$ 锁定于恒星中时，风带走
$\lambda\,\mathrm{d}s$ 的气体，\omterm{def:b3:origin-of-elements:metallicity}{金属丰度}为 $Z$：
$\mathrm{d}M_\mathrm{g} = -(1 + \lambda)\,\mathrm{d}s$，且
$\mathrm{d}(ZM_\mathrm{g}) = y\,\mathrm{d}s - Z\,\mathrm{d}s -
\lambda Z\,\mathrm{d}s$。将第二式展开，再
减去 $Z\,\mathrm{d}M_\mathrm{g}$，剩下 $M_\mathrm{g}\,\mathrm{d}Z =
y\,\mathrm{d}s = -\frac{y}{1 + \lambda}\,\mathrm{d}M_\mathrm{g}$，积分后
即得所求结果。(b)~$\ln(1/0.29) = 1.238$：$Z = 0.0155$ 对应
$\lambda = 0$，减半后为 $0.0077$，对应 $\lambda = 1$。(c)~风
会使有效产额降低 $1 + \lambda$ 倍。在质量
较小的星系中，超新星更易将气体驱出，$\lambda$ 较大，
金属随气体一并流失。
\end{solution}

\begin{solution}{exo:b3:origin-of-elements:12}
(a)~太阳中 $N_\mathrm{Li}/N_\mathrm{H} = 10^{0.96 - 12} = \num{9.1e-12}$，
CI 球粒陨石中为 $10^{3.25 - 12} = \num{1.78e-9}$。太阳
保留了其锂的 $10^{-2.29} = \num{0.0051}$，即已破坏了
\qty{99.5}{\percent} 的锂。(b)~$\num{1.78e-9}/\num{1.45e-10} = 12.3$，以及
$\num{1.78e-9}/\num{4.72e-10} = 3.8$。(c)~太阳系所含的锂
几乎是大爆炸所造锂量的四倍：其锂的大部分
是之后由其他过程产生的。最古老的恒星所含锂
比大爆炸所造的少三倍：它们或如太阳那样破坏了其中一部分，
或是预测遗漏了某些因素；仅凭这两个比较
无法断定究竟是哪一种。
\end{solution}

\begin{solution}{pb:b3:origin-of-elements:1}

\textbf{第一部分。}

\textbf{1.} $10^{12 - 7.46} = 10^{4.54} = \num{34700}$.

\textbf{2.} $N_\mathrm{He}/N_\mathrm{H} = 10^{10.914 - 12} = \num{0.0820}$，于是
$4N_\mathrm{He}/N_\mathrm{H} = \num{0.328}$，对比 $Y/X =
0.2423/0.7438 = \num{0.326}$。二者在 \qty{0.7}{\percent} 以内一致；其间的
微小差距之所以存在，是因为一个氦原子的质量略小于
四个氢原子。

\textbf{3.} $\log_{10}\bigl((1/350000)/\num{2.884e-5}\bigr) = -5.544 + 4.540 =
-1.00$.

\textbf{4.} $[\mathrm{O/H}] = -1.00 + 0.40 = -0.60$，故
$\log\varepsilon(\mathrm{O}) = 8.69 - 0.60 = 8.09$。该恒星含有
每个铁原子对应 $10^{8.09 - 6.46} = 42.7$ 个氧原子，而
在太阳中为 $10^{1.23} = 17.0$，该恒星的比值是太阳的 $10^{0.4} = 2.5$ 倍。

\textbf{5.} $Z \approx 0.0139\times10^{-0.60} = \num{0.0035}$，约为太阳的
四分之一。

\textbf{6.} 在低温恒星的光谱中，铁有非常多的谱线，而
氧的谱线少且弱：铁的测量要容易得多，
也更精确，因此可代表\omterm{def:b3:origin-of-elements:metallicity}{金属丰度}。

\textbf{第二部分。}

\textbf{7.} 每个核都位于一条较低 $Z$ 的稳定等压线之后：
$^{122}\mathrm{Sn}$、$^{123}\mathrm{Sb}$ 和 $^{124}\mathrm{Sn}$ 都会
使 \omterm{def:b3:origin-of-elements:s-process}{r 过程}的衰变链在到达碲之前中断。三者都位于
s 路径上，该路径自 $^{122}\mathrm{Te}$ 起经过碲，直至
$^{126}\mathrm{Te}$。

\textbf{8.} \num{31.48}、\num{31.03} 和 \qty{30.75}{mb}；平均值为
\qty{31.09}{mb}.

\textbf{9.} 另外两者的平均值为 $(31.48 + 30.75)/2 = \qty{31.11}{mb}$，
故 $N_{123} = 31.11/832 = \num{0.03740}$，仅比
测得值 \num{0.0373} 高 \qty{0.3}{\percent}。

\textbf{10.} $1.329\times44.4 = \qty{59.0}{mb}$，几乎是平台值的两倍。
$^{128}\mathrm{Te}$ 仅由 \omterm{def:b3:origin-of-elements:s-process}{r 过程}产生：s 路径在
$^{127}\mathrm{Te}$ 处离开碲，因此其丰度由 \omterm{def:b3:origin-of-elements:s-process}{r 过程}产生，
与 s 流无关。

\textbf{11.} $1 - 7.43/8.53 = \num{0.13}$：\qty{13}{\percent} 的流
绕过了 $^{148}\mathrm{Sm}$，经由对不稳定核
$^{147}\mathrm{Nd}$、$^{147}\mathrm{Pm}$ 与 $^{148}\mathrm{Pm}$ 的俘获，
这些核都来不及衰变为钐。

\textbf{12.} 中子数为 82 的闭壳层使这些核
只较弱地束缚额外的一个中子，且俘获截面很小（\qty{4.0}{mb}
对于 $^{138}\mathrm{Ba}$ 为）。中子流在它们上面堆积，
只有一部分继续传递：它们是瓶颈。

\textbf{第三部分。}

\textbf{13.} $\mu_1 = 13.7/(13.7 + 33.4) = 13.7/47.1 = \num{0.291}$.

\textbf{14.} $\ln(1/0.291) = \num{1.235}$，故 $y = 0.0154/1.235 = \num{0.0125}$，
为 $Z_\odot$ 的 \num{0.81} 倍。

\textbf{15.} 当 $\mu = 0.9$ 时，$\ln(1/0.9) = 0.1054$，故 $Z = 0.0125\times0.1054 =
\num{0.00131}$。由 $0.00131/0.0154 = 0.0851$，得 $[\mathrm{Fe/H}] =
-1.07$。当 $\mu = 0.5$ 时，$Z = 0.0125\times0.693 = \num{0.00864}$，
则 $[\mathrm{Fe/H}] = -0.25$。

\textbf{16.} $Z = 0.00154$，故 $\mu = \mathrm{e}^{-0.00154/0.01247} =
\mathrm{e}^{-0.1235} = \num{0.884}$。

\textbf{17.} $1 - 0.884 = \num{0.116}$，即质量的 \qty{12}{\percent}。

\textbf{18.} $y = 0.0154/\ln 10 = \num{0.0067}$，约为前者的一半。

\textbf{第四部分。}

\textbf{19.} \omterm{def:b3:origin-of-elements:metallicity}{金属丰度}低于 $Z$ 的恒星质量为
$M(1 - \mathrm{e}^{-Z/y})$，故单位 $Z$ 对应的质量为
$(M/y)\,\mathrm{e}^{-Z/y}$。由 $Z = Z_\odot10^{[\mathrm{Fe/H}]}$ 可得
$\mathrm{d}Z = Z\ln 10\,\mathrm{d}[\mathrm{Fe/H}]$，故单位
$[\mathrm{Fe/H}]$ 的分布正比于
$Z\mathrm{e}^{-Z/y}$。其导数
$(1 - Z/y)\mathrm{e}^{-Z/y}$ 在 $Z = y$ 处为零，峰值位于
$[\mathrm{Fe/H}] = \log_{10}0.810 = -0.09$。

\textbf{20.} $\num{0.116}/(1 - 0.291) = 0.1162/0.7091 = \num{0.164}$。

\textbf{21.} $Z/y = 1.235\times10^{-0.5} = 0.3905$，$1 - \mathrm{e}^{-0.3905} =
0.3233$，且 $0.3233/0.7091 = \num{0.456}$。

\textbf{22.} $6670\times0.1638 = 1093$，与 26 相比，多了约四十倍。

\textbf{23.} $26/902 = \qty{2.9}{\percent}$，仍比
方框中的值低五倍以上。模型中的恒星都是长寿命的，即质量至多约一个太阳质量的那些，
它们自形成以来均已存活至今；在质量较大的 F 型星中，
最古老者（也是\omterm{def:b3:origin-of-elements:metallicity}{金属丰度}最低者）
已经死亡，完整样本中它们的数目被低估了。

\textbf{24.} $(1 - 0.1^{0.1})/(1 - 0.1) = 0.2057/0.9 = \num{0.23}$，更差。
提高气体比例也无济于事：当 $\mu_1$ 趋于 1 时，
该比例趋于 $0.1$，仍为
观测值 \qty{2.9}{\percent} 的三倍以上。任何气体比例都无法挽救封\omterm{thm:b3:origin-of-elements:closed-box}{闭箱模型}。

\textbf{25.} 封\omterm{thm:b3:origin-of-elements:closed-box}{闭箱模型}预言，
\textbf{\qty[text-series-to-math=true]{16}{\percent}} 的长寿命
恒星中 $[\mathrm{Fe/H}] < -1$ 者占，而观测到的至多
为 \qty{3}{\percent}。盘面并非封闭箱：它
通过缓慢吸积贫金属气体而增长，因此很少有气体
曾停留在低\omterm{def:b3:origin-of-elements:metallicity}{金属丰度}状态下等待成星。
\end{solution}
